
\documentclass[11pt]{article}


\usepackage[utf8]{inputenc}
\usepackage{textcomp}

\usepackage{url}
\usepackage{hyperref}
\usepackage{graphicx}
\usepackage{grffile}


\graphicspath{{figs/}}

\usepackage{url}
\usepackage{hyperref}

\date{2017-Jun-25}
\title{Study of Manuscript Creating Using Latex Word Processing }
\usepackage{authblk}
\author[1]{Author A\thanks{A.A@university.edu}}
\author[1]{Author B\thanks{B.B@university.edu}}
\author[1]{Author C\thanks{C.C@university.edu}}
\author[2]{Author D\thanks{D.D@university.edu}}
\author[2]{Author E\thanks{E.E@university.edu}}
\affil[1]{Department of Computer Science, \LaTeX\ University}
\affil[2]{Department of Mechanical Engineering, \LaTeX\ University}

\begin{document}

\maketitle

\begin{abstract}
In this work the method to create a manuscript will be demonstrated.
 The use of Unibeam, combined with pdflatex and bibtex, under Linux will be used to create a manuscript. 
\end{abstract}


\section{Introduction}
This is the introduction of the manuscript.
The first citation will be given here.\cite{spartrekus2324898}
There second citation will be here.\cite{spartrekus2324897}

\section{Description}
There are several templates for Ph.d. thesis, which are available in the 
reference Spartrekus~\cite{spartrekus2324898}.
It describes 11 examples of Ph.d. thesis templates. The first templates 
are very basic and it allows make a less complex manuscript.
The example 05 gives a cover with a picture of the university.
The examples, numbered with more than ex.6, are more complex and 
they give a beautiful looking.

The most difficult part is probably the references which need to be 
formatted in order appearance. 
There is a nice example from Spartrekus~\cite{spartrekus2324897}.
This example gives a Makefile, which allows to automatically bring
the manuscript in form, and to display it using mupdf. 

\begin{center}
\includegraphics[width=1.0\textwidth]{freesoft.jpg}
\end{center}

\section{Conclusion}
It allowed to easily create a manuscript and to make use of a bib file, which contains references.

\bibliography{mybib}{}
\bibliographystyle{ieeetr}

\end{document}
 

